% Part: model-theory
% Chapter: models-of-arithmetic
% Section: computable-models

\documentclass[../../../include/open-logic-section]{subfiles}

\begin{document}

\olfileid{mod}{mar}{cmp}
\section{అంకగణితపు గణనీయ నమూనాలు}

\begin{explain}
ప్రామాణిక నమూనా~$\Struct{N}$కు రెండు మంచి లక్షణాలు ఉన్నాయి. దాని
వ్యక్తి క్షేత్రం సహజ సంఖ్యలు~$\Nat$; అంటే అంకగణిత భాషతో మనం ఏ రకపు
వస్తువులను గురించి మాట్లాడాలనుకుంటామో దాని మూలకాలు సరిగ్గా అవే,
మరియు ప్రామాణిక సంఖ్యాపదం $\num{n}$ నిజంగానే $n$ను ఎంచుతుంది.
మరొక మంచి లక్షణం: $\Lang{L_A}$ తార్కికేతర సంకేతాల అర్థనిర్దేశాలన్నీ
\emph{గణనీయాలు}. $\Assign{\prime}{N}$, $\Assign{+}{N}$,
$\Assign{\times}{N}$లుగా పనిచేసే ఉత్తరాధికారి, సంకలన, గుణకార
ప్రమేయాలు సంఖ్యల గణనీయ ప్రమేయాలు. (ఉదాహరణకు, ట్యూరింగ్ యంత్రాలతో
గణించదగినవి లేదా ఆదిమ పునరావర్తనంతో నిర్వచించదగినవి.) అలాగే
$\Struct{N}$పై తక్కువ-కంటే సంబంధం, అంటే $\Assign{<}{N}$,
నిర్ణయించదగినది.
\end{explain}

$\Th{Q}$, $\Th{PA}$ వంటి అంకగణిత సిద్ధాంతాల అప్రామాణిక నమూనాలు
అప్రామాణిక మూలకాలను తప్పక కలిగి ఉండాలి. అందువల్ల వాటి వ్యక్తి
క్షేత్రాల్లో సాధారణంగా $\Nat$కు అదనంగా
\tetoken{మూలకాలు}{!!{element}s} ఉంటాయి. అయితే ప్రతి లెక్కించదగిన
నిర్మాణాన్ని $\Nat$తో సహా ఏ
\tetoken{అనంతంగా లెక్కించదగిన}{!!{denumerable}} సమితిపైనైనా
నిర్మించవచ్చు. కాబట్టి వ్యక్తి క్షేత్రం $\Nat$గానే ఉన్న
అప్రామాణిక నమూనాలు కూడా ఉన్నాయి. అలాంటి $\Struct{M}$ నమూనాల్లో
కనీసం కొన్ని సంఖ్యలు తమ సాధారణ పాత్రలను పోషించలేవు; ఎందుకంటే ఏదో
$k$ ఉండాలి, అది ప్రతి $n\in\Nat$కూ $\Value{\num{n}}{M}$ కంటే
భిన్నంగా ఉండాలి.

\begin{defn}
$\Lang{L_A}$కు గల
\tetoken{ఒక నిర్మాణం}{!!^a{structure}}~$\Struct{M}$కు
$\Domain{M}=\Nat$ అయి, $\Assign{\prime}{M}$, $\Assign{+}{M}$,
$\Assign{\times}{M}$లు గణనీయ ప్రమేయాలు, $\Assign{<}{M}$
నిర్ణయించదగిన సంబంధం అయితే, అప్పుడూ అప్పుడే అది \emph{గణనీయం}.
\end{defn}

\begin{ex}\ollabel{ex:comp-model-q}
\olref[mdq]{ex:model-K-of-Q}లోని
\tetoken{నిర్మాణం}{!!{structure}}~$\Struct{K}$ను గుర్తుచేసుకోండి.
దాని వ్యక్తి క్షేత్రం $\Domain{K}=\Nat\cup\{a\}$; అర్థనిర్దేశాలు
\begin{align*}
  \Assign{\Obj{0}}{K} & = 0\\
  \Assign{\prime}{K}(x) & =
  \begin{cases}
    x+1 & \text{$x\in\Nat$ అయితే}\\
    a & \text{$x=a$ అయితే}
  \end{cases}\\
  \Assign{+}{K}(x, y) & =
  \begin{cases}
    x+y & \text{$x,y\in\Nat$ అయితే}\\
    a & \text{మిగతా సందర్భాల్లో}
  \end{cases}\\
  \Assign{\times}{K}(x, y) & =
  \begin{cases}
    xy & \text{$x,y\in\Nat$ అయితే}\\
    0 & \text{$x=0$ లేదా $y=0$ అయితే}\\
    a & \text{మిగతా సందర్భాల్లో}\\
  \end{cases}\\
  \Assign{<}{K} & =
  \Setabs{\tuple{x,y}}{x, y \in \Nat \text{ మరియు } x<y} \cup
  \Setabs{\tuple{x,a}}{x \in \Domain{K}}
\end{align*}
\sourcecorrection{OLTEMODARI-013}{చివరి సమితి-నిర్మాణంలో క్రమయుగ్మం
మొదటి ఘటకం $x$ అయినప్పటికీ ఆంగ్ల మూల షరతు సంబంధం లేని
$n\in\Domain{K}$ను వాడుతుంది. ఇదే నిర్మాణాన్ని ముందుగా ఇచ్చిన సరైన
నిర్వచనానికి అనుగుణంగా $x\in\Domain{K}$గా సరిచేశాం.}
అయితే $\Domain{K}$
\tetoken{అనంతంగా లెక్కించదగినది}{!!{denumerable}}; కాబట్టి అది
$\Nat$తో సమసంఖ్యకం. ఉదాహరణకు, $g(0)=a$, $n>0$కు $g(n)=n-1$
అయ్యే $g\colon\Nat\to\Domain{K}$ ఒక
\tetoken{ద్విగుణ ప్రమేయం}{!!a{bijection}}. దీన్ని ఉపయోగించి
$\Th{Q}$ యొక్క కొత్త నమూనా~$\Struct{K'}$కూ $\Struct{K}$కూ మధ్య
సమరూపతను ఏర్పరచవచ్చు. $\Struct{K'}$లో $\Lang{L_A}$ సంకేతాలకు వేరే
ప్రమేయాలు, సంబంధాలు కేటాయించాలి; ఎందుకంటే $\Nat$లోని వేరే
\tetoken{మూలకాలు}{!!{element}s} ప్రామాణిక, అప్రామాణిక సంఖ్యల
పాత్రలను పోషిస్తాయి.
\sourcecorrection{OLTEMODARI-015}{ఆంగ్ల మూలంలోని
$g(0)=a$, $g(n)=n+1$ చిత్రణ $0$, $1$లను తాకదు; అందువల్ల అది
ద్విగుణ ప్రమేయం కాదు, తరువాత ఇచ్చిన తరలించిన క్రియలనూ ఉత్పత్తి చేయదు.
$g(n)=n-1$ (అన్ని $n>0$లకు) మాత్రమే $\Nat$ నుంచి
$\Nat\cup\{a\}$కు ద్విగుణ ప్రమేయమై, కిందటి $x+y-1$ మరియు
$xy-x-y+2$ సూత్రాలను ఇస్తుంది; ఆ ఉద్దేశిత చిత్రణను వాడాం.}

మరింత నిర్దిష్టంగా, ఇప్పుడు $0$ అతి చిన్న ప్రామాణిక సంఖ్య పాత్రను
కాక $a$ పాత్రను పోషిస్తుంది. అతి చిన్న ప్రామాణిక సంఖ్య ఇప్పుడు
$1$. అందువల్ల $\Assign{\Obj{0}}{K'}=1$గా కేటాయిస్తాం.
ఉత్తరాధికారి ప్రమేయం కూడా ఇప్పుడు వేరు: ప్రామాణిక సంఖ్య, అంటే
$n>0$, ఇస్తే అది ఇంకా $n+1$నే ఇస్తుంది. కానీ ఇప్పుడు $0$ తనకే
ఉత్తరాధికారి అయిన $a$ పాత్రను పోషిస్తుంది; కాబట్టి
$\Assign{\prime}{K'}(0)=0$. సంకలనం, గుణకారాలకు కూడా
\begin{align*}
\Assign{+}{K'}(x, y) & =
  \begin{cases}
    x+y-1 & \text{$x,y>0$ అయితే}\\
    0 & \text{మిగతా సందర్భాల్లో}
  \end{cases}\\
  \Assign{\times}{K'}(x, y) & =
  \begin{cases}
    1 & \text{$x=1$ లేదా $y=1$ అయితే}\\
    xy - x - y + 2 & \text{$x,y>1$ అయితే}\\
    0 & \text{మిగతా సందర్భాల్లో}\\
  \end{cases}
\end{align*}
అలాగే $\tuple{x,y}\in\Assign{<}{K'}$ అవడం సరిగ్గా
$x<y$, $x>0$, $y>0$ అన్నీ నిజమైనప్పుడు, లేదా $y=0$ అయినప్పుడు.

ఈ ప్రమేయాలన్నీ సహజ సంఖ్యల గణనీయ ప్రమేయాలు; $\Assign{<}{K'}$,
$\Nat$పై నిర్ణయించదగిన సంబంధం. కానీ అవి $\Nat$పై ఉత్తరాధికారి,
సంకలన, గుణకార ప్రమేయాలు కావు; అలాగే $\Assign{<}{K'}$,
$\Nat$పైని $<$ సంబంధం కాదు.
\end{ex}

\begin{prob}
\olref[mod][mar][mdq]{ex:model-L-of-Q}లోని $\Struct{L}$తో సమరూపమై,
$\Domain{L'}=\Nat$ అయ్యే ఒక
\tetoken{నిర్మాణం}{!!a{structure}}~$\Struct{L'}$ ఇవ్వండి.
\end{prob}

\begin{explain}
$\Th{Q}$కు వ్యక్తి క్షేత్రం $\Nat$గానే ఉన్న గణనీయ అప్రామాణిక
నమూనాలు ఉన్నాయని \Olref{ex:comp-model-q} చూపుతుంది. అయితే
$\Th{PA}$ నమూనాలకు (అందువల్ల $\Th{TA}$ నమూనాలకు కూడా) ఇది నిజం
కాదని కింది ఫలితం చూపుతుంది.
\end{explain}

\begin{thm}[టెన్నెన్‌బామ్ సిద్ధాంతం]
$\Th{PA}$ యొక్క ప్రతి గణనీయ నమూనా $\Struct{N}$తో సమరూపమైన
ప్రామాణిక నమూనా.
\sourcecorrection{OLTEMODARI-014}{ఆంగ్ల మూలం $\Struct{N}$ మాత్రమే
గణనీయ నమూనా అంటుంది. కానీ $\Nat$పై ఒక గణనీయ ద్విగుణ ప్రమేయం ద్వారా
$\Struct{N}$ నిర్మాణాన్ని తరలిస్తే, దానితో సమరూపమైనప్పటికీ అక్షరాలా
అదే కాని గణనీయ ప్రతులు వస్తాయి. టెన్నెన్‌బామ్ సిద్ధాంతం తొలగించేది
గణనీయ \emph{అప్రామాణిక} నమూనాలను; అందుకే ఖచ్చితమైన
“ప్రతి గణనీయ నమూనా ప్రామాణికం” రూపాన్ని ఇచ్చాం.}
\end{thm}

\end{document}
